\section{A Composite Score Rejects Costly Gains}

$Q_{\mathrm{HW}}$ is the implementation's bounded promotion utility, not a
hardware-only objective. For anchor $a$ and candidate $c$, let
$r_L=(p_aL_a)/(p_cL_c)$ use cycle latency $L$ and effective clock period $p$.
When initiation interval applies,
$r_P=r_L^{0.85}(II_a/II_c)^{0.15}$; otherwise $r_P=r_L$. Define the
capacity-normalized resource footprint
$A(x)=\sum_r x_r/K_r$ over LUT, FF, DSP, BRAM, and URAM, and
$r_A=A(a)/A(c)$. Let $D\in\{0,1\}$ mark a source change and
$R\in\{0,1\}$ mark a valid candidate that repairs an invalid starter. The
composite ratio is
$h=0.99^D2^Rr_P^{0.55}r_A^{0.45}$.

$Q_{\mathrm{HW}}=1-1/(1+h)^2$ maps $h{=}1$ to $0.75$, approaches $1$ for
large gains, and approaches $0$ for regressions. Budget efficiency
$E=\max(0.80,1-0.10u_{\mathrm{credit}}-0.10u_{\mathrm{time}})$ uses consumed
fractions clipped to $[0,1]$. The final score is
\begin{equation}
S=100\,V(c)\,Q_{\mathrm{HW}}\,E.
\label{eq:score}
\end{equation}
The 0.99 and 2 factors encode anti-no-op and repair policies, while 0.85/0.15
is a fixed latency/II split. The 0.55/0.45 performance/resource balance was
heuristically calibrated against 95 official PPA references using 99 fresh
Vitis~2025.2 measurements; 0.55 lies inside their feasible interval. This
supports a common operating point, but the evidence is limited and one fixed
exponent need not fit every task. Promotion occurs only after $V(c){=}1$, and
a failed gate yields $S{=}0$.
